% !TeX root = ../../Book5.tex
% 纲要标题：### 11.4 结构问题
% 纲要位置：工作记录/计划纲要.md，第 4105 行。

\section{结构问题}
\label{b5:sec:1104}


% 纲要范围：
%
% * 可解、幂零及半单的定义
% * 导出代数与导出列
% * 复 Lie 代数与一般正特征 Lie 代数的区别
%
% ---
%

最初的结构问题，是反复取括号以后会发生什么。导出列让两边同时来自上一步，下降中心列则始终允许一边来自整个代数。二者给出的有限终止条件不同；弄清这种区别，才能准确使用下一章的三角化定理。

\subsection{可解、幂零与半单 Lie 代数}
\label{b5:subsec:1104:01}

\begin{definition}\label{b5:def:lie-solvable-nilpotent}
设 \(\mathfrak g\) 为域 \(k\) 上 Lie 代数。定义导出列
\[
\mathfrak g^{(0)}=\mathfrak g,\qquad
\mathfrak g^{(r+1)}=\bigl[\mathfrak g^{(r)},\mathfrak g^{(r)}\bigr],
\]
以及下降中心列
\[
\gamma_1(\mathfrak g)=\mathfrak g,\qquad
\gamma_{r+1}(\mathfrak g)=\bigl[\mathfrak g,\gamma_r(\mathfrak g)\bigr].
\]
若存在整数 \(r\ge0\) 使 \(\mathfrak g^{(r)}=0\)，称 \(\mathfrak g\) \term[kejie-lie-daishu]{可解}[solvable]。此时记
\[
\ell(\mathfrak g)=\min\{r\ge0:\mathfrak g^{(r)}=0\},
\]
称它为 \(\mathfrak g\) 的\term[kejie-changdu]{可解长度}[derived length]。若存在整数 \(c\ge0\) 使 \(\gamma_{c+1}(\mathfrak g)=0\)，称 \(\mathfrak g\) \term[miling-lie-daishu]{幂零}[nilpotent]。此时最小的这种 \(c\) 记为 \(c(\mathfrak g)\)，称为 \(\mathfrak g\) 的幂零类。零 Lie 代数的可解长度与幂零类均为 \(0\)。
\end{definition}
\symidx{\(\mathfrak g^{(r)},\gamma_r(\mathfrak g)\)：Lie 代数的导出列与下降中心列}
\symidx{\(\ell(\mathfrak g),c(\mathfrak g)\)：Lie 代数的可解长度与幂零类}
\begin{definition}\label{b5:def:lie-semisimple-simple}
设 \(\mathfrak g\) 为域 \(k\) 上有限维 Lie 代数。若它没有非零可解理想，就称为\term[bandan-lie-daishu]{半单 Lie 代数}[semisimple Lie algebra]。若它非交换且只有零与自身两个理想，就称为\term[dan-lie-daishu]{单 Lie 代数}[simple Lie algebra]。
\end{definition}
这里把非交换性写入“单”的定义，因而一维交换 Lie 代数不称为单 Lie 代数。半单代数最终能否分解为单理想之和，还需要证明；在定义中先使用理想条件，避免把后面的分解定理预先假定下来。
\ProseParagraphBreak
三维 Heisenberg 代数满足
\(\gamma_2=kz,\gamma_3=0\)，所以幂零类为二。二维代数 \([h,e]=e\) 则满足
\[
\mathfrak g^{(1)}=ke,\quad \mathfrak g^{(2)}=0,\qquad
\gamma_r(\mathfrak g)=ke\quad(r\ge2).
\]
它可解但不幂零。这里的一维理想 \(ke\) 和一维商 \(\mathfrak g/ke\) 都是交换 Lie 代数，已经说明幂零性不能像可解性一样在任意扩张下保持。

\subsection{导出代数与导出列}
\label{b5:subsec:1104:02}

\begin{proposition}\label{b5:prop:lie-derived-ideal}
设 \(\mathfrak g\) 为域 \(k\) 上的 Lie 代数。\([\mathfrak g,\mathfrak g]\) 是理想，称为 \(\mathfrak g\) 的\term[daochu-daishu]{导出代数}[derived algebra]；商
\(\mathfrak g/[\mathfrak g,\mathfrak g]\)
交换；从 \(\mathfrak g\) 到任意交换 Lie 代数的同态都唯一通过这个商分解。
\end{proposition}
\begin{proof}
Jacobi 的导子形式说明
\(\bigl[z,[x,y]\bigr]=\bigl[[z,x],y\bigr]+\bigl[x,[z,y]\bigr]\)
仍在括号的线性生成空间中，所以它为理想。商中全部括号为零。到交换代数的同态把每个括号送到零，核包含该理想，商映射的泛性质即给出唯一分解。
\end{proof}
\begin{proposition}\label{b5:prop:lie-series-properties}
设 \(\mathfrak g\) 为域 \(k\) 上的 Lie 代数，记 \(\mathfrak g^{(r)}\) 为导出列的第 \(r\) 项，\(\gamma_r(\mathfrak g)\) 为下降中心列的第 \(r\) 项。两列的每一项都是理想；对整数 \(r,s\ge1\)，有
\[
\bigl[\gamma_r(\mathfrak g),\gamma_s(\mathfrak g)\bigr]
\subseteq\gamma_{r+s}(\mathfrak g),\qquad
\mathfrak g^{(r)}\subseteq\gamma_{2^r}(\mathfrak g).
\]
特别地，幂零蕴涵可解。子代数与商分别继承可解性及幂零性。
\end{proposition}
\begin{proof}
若 \(I,J\) 为 \(\mathfrak g\) 的理想，Jacobi 恒等式给出
\[
\bigl[\mathfrak g,[I,J]\bigr]
\subseteq\bigl[[\mathfrak g,I],J\bigr]+\bigl[I,[\mathfrak g,J]\bigr]
\subseteq[I,J],
\]
故 \([I,J]\) 仍为理想。从 \(\mathfrak g\) 自身是理想出发，分别按导出列和下降中心列的递推定义归纳，得到两列的每一项都是理想。

对第一个包含关系按 \(r\) 归纳，\(r=1\) 是定义；由 Jacobi，
\[
\bigl[[\mathfrak g,\gamma_r],\gamma_s\bigr]
\subseteq\bigl[\mathfrak g,[\gamma_r,\gamma_s]\bigr]
+\bigl[\gamma_r,[\mathfrak g,\gamma_s]\bigr]
\subseteq\gamma_{r+s+1}.
\]
第二项使用同一归纳假设对所有 \(s\) 成立。随后按导出次数归纳：
若 \(\mathfrak g^{(r)}\subseteq\gamma_{2^r}\)，两个这样的项取括号便落在 \(\gamma_{2^{r+1}}\)。若下降中心列终止，导出列因此也终止。
子代数的两种列逐项包含于原代数对应的列；满同态则将两种列逐项映满商中的相应列。故终止性质在这两种操作下保持。
\end{proof}
两种列的相邻商说明了它们不同的结构含义。导出列满足
\[
[\mathfrak g^{(r)},\mathfrak g^{(r)}]
=\mathfrak g^{(r+1)},
\]
所以 \(\mathfrak g^{(r)}/\mathfrak g^{(r+1)}\) 是交换 Lie 代数。下降中心列则满足 \([\mathfrak g,\gamma_r(\mathfrak g)]=\gamma_{r+1}(\mathfrak g)\)，因而
\[
\gamma_r(\mathfrak g)/\gamma_{r+1}(\mathfrak g)
\subseteq Z\bigl(\mathfrak g/\gamma_{r+1}(\mathfrak g)\bigr).
\]
后者不仅要求这一层内部括号为零，还要求它在相应商中与整个代数括号为零。这也解释了为什么仅由交换层组成的扩张可以可解而不幂零。

\begin{proposition}\label{b5:prop:lie-solvable-extension}
设 \(\mathfrak g\) 为域 \(k\) 上的 Lie 代数，\(\mathfrak a\) 为其理想。若 \(\mathfrak a\) 与 \(\mathfrak g/\mathfrak a\) 都可解，则 \(\mathfrak g\) 可解。若 \(\mathfrak a\subseteq Z(\mathfrak g)\)，且 \(\mathfrak g/\mathfrak a\) 幂零，则 \(\mathfrak g\) 幂零。
\end{proposition}
\begin{proof}
若商的第 \(r\) 个导出项为零，则 \(\mathfrak g^{(r)}\subseteq\mathfrak a\)。再取 \(s\) 次导出，落入 \(\mathfrak a^{(s)}\)，最终为零。对第二个结论，若商的第 \(c+1\) 个下降中心项为零，则
\(\gamma_{c+1}(\mathfrak g)\subseteq\mathfrak a\subseteq Z(\mathfrak g)\)，从而 \(\gamma_{c+2}(\mathfrak g)=0\)。
\end{proof}
\begin{proposition}\label{b5:prop:lie-simple-perfect}
设 \(\mathfrak g\) 为域 \(k\) 上的非交换有限维单 Lie 代数。则
\([\mathfrak g,\mathfrak g]=\mathfrak g\)、\(Z(\mathfrak g)=0\)，且 \(\mathfrak g\) 半单。域 \(k\) 上任意有限维半单 Lie 代数的中心也为零。
\end{proposition}
\begin{proof}
导出代数为非零理想，所以等于整个 \(\mathfrak g\)。其导出列不终止，故 \(\mathfrak g\) 不可解；唯一可能的非零理想因此不是可解理想。中心为理想，若等于整个代数便使代数交换，所以只能为零。一般半单代数的中心是交换理想，同样必须为零。
\end{proof}

\subsection{复数域与正特征的区别}
\label{b5:subsec:1104:03}

本编关于半单结构与有限维表示的核心定理以复数域为环境。特征零使某些整数系数可以相消，代数闭性使特征值存在；两种条件有不同用途。
\ProseParagraphBreak
例如在特征 \(p>0\) 时，\(\mathfrak{sl}_p(k)\) 包含非零标量矩阵 \(I_p\)，因为 \(\operatorname{tr}I_p=p=0\)。它张成一个中心交换理想，所以这个 Lie 代数不满足本节的半单定义。尤其在特征二的 \(\mathfrak{sl}_2(k)\) 中，取
\[
e=E_{12},\qquad f=E_{21},\qquad h=I_2,
\]
便有 \([e,f]=h\)、\([h,e]=[h,f]=0\)。它成为三维 Heisenberg 代数，而不是后面要研究的复 \(\mathfrak{sl}_2\)。
\ProseParagraphBreak
还要区分“代数幂零”与“在某个表示中每个算子幂零”。一个一维交换 Lie 代数本身幂零，却可让它的基向量在一维表示上作用为一，得到非幂零算子。Engel 定理将要求给定表示中的每个作用算子幂零；应用于伴随表示以后，才能得到 Lie 代数自身的幂零性判据。下一章会把这个转换写清楚。
